\cleardoublepage
\chapternonum{摘要}

自动驾驶技术是人工智能和机器人领域的研究热点之一，而自动驾驶赛车是验证高阶自动驾驶技术的重要平台。自动驾驶赛车聚焦高速极限工况下的运动控制问题，能够推动自动驾驶系统在安全性、稳定性和性能方面的技术进步。然而，赛车在高速运动时，具有非线性的动力学特性和复杂的环境交互，导致运动控制面临诸多挑战。尤其是在未知复杂环境中，如何实现安全性与竞速性能的协同优化成为一个亟待解决的核心问题。本文针对这一问题，系统研究了从已知赛道到未知复杂环境、由竞速控制到高速导航的多层次运动控制方法。本文的主要研究内容如下：

\begin{enumerate}
    \item 针对已知动态环境中的安全竞速问题，本文首先建立了赛车的动力学模型，包括车身动力学模型和轮胎力学模型，并通过最小二乘辨识方法对模型参数进行估计和优化。基于该模型，本文提出了基于控制障碍函数的安全竞速控制方法，构建安全集并从理论上保证赛车在高速运动中的安全性。同时，本文进一步引入了基于控制障碍函数的模型预测控制方法，将安全约束融入优化问题中，通过引入进度变量并最大化赛道进度来实现安全的高速竞速。实验在仿真和实车环境中验证了本文方法，结果表明该方法能够显著提升赛车的安全性，同时保持较高的竞速性能。
    \item 针对未知封闭赛道中的视觉竞速问题，本文提出了一种两阶段的基于强化学习的端到端竞速方法。在第一阶段，本文使用强化学习算法训练了一个教师策略，该策略利用仿真中的特权信息实现了高性能的竞速控制。在第二阶段，本文采用知识蒸馏的方法将教师策略迁移到一个仅依靠视觉的学生策略，使其能在未知的赛道中实现竞速。本文在提升视觉控制策略的训练效率的同时，解决了部分可观环境中难以获得最优控制的难题。为了提高学生策略的泛化能力，本文还引入了域随机化技术，在训练过程中对环境参数进行随机化。实验结果表明，所提出的方法在仿真环境中表现出优于传统控制方法和其他端到端方法的竞速性能，并且成功实现了仿真到现实的零样本迁移。
    \item 针对未知复杂环境中的点到点高速导航问题，本文提出了一种基于视觉预测模型的端到端高速导航方法。首先，本文设计了一个循环状态空间模型来预测未来的环境状态。该方法通过收集轨迹数据，利用视觉输入学习环境的动态特性，从而提升预测的准确性。基于该视觉预测模型，本文引入了模型预测路径积分控制方法，通过采样动作空间获得预测轨迹来优化控制策略。此外，针对高速导航问题，本文设计了一个奖励函数，用于评估采样轨迹的质量，从而指导控制策略的优化。该方法解决了强化学习方法的训练效率问题，以及部分可观环境中长距离导航的短视问题。最后，本文通过对比与消融实验，凸显了方法的优越性，并且在实车环境中验证了其有效性。

\end{enumerate}

\noindent \textbf{关键词}：自动驾驶赛车；模型预测控制；控制障碍函数；强化学习；端到端导航

\cleardoublepage
\chapternonum{Abstract}

Autonomous driving technology is an important research topic in the fields of artificial intelligence and robotics, and autonomous racing cars have become an important platform for validating advanced autonomous driving technologies. Autonomous racing research focuses on the motion control problem under high-speed extreme conditions, which can promote technological progress in the safety, stability, and performance of autonomous driving systems. However, racing cars have highly nonlinear dynamics and complex interactions with the environment during high-speed motion, leading to many challenges in motion control. In particular, how to balance safety and racing performance in unknown complex environments has become a core problem that warrants further exploration. This thesis systematically studies multi-level motion control methods from known tracks to unknown complex environments, from racing control to high-speed navigation, and proposes a series of algorithms and techniques. The main research contents of this thesis are as follows:

\begin{enumerate}
    \item For the safe racing control in known dynamic environments, this thesis first establishes a dynamic model of the racing car, including a vehicle dynamics model and a tire dynamics model, and optimizes the model parameters using a least squares identification method. Based on this model, this thesis proposes a control barrier function method to construct a safe set and theoretically guarantee the safety of the high-speed racing car. Meanwhile, this thesis introduces a model predictive control method based on control barrier functions, which incorporates safety constraints into the optimization problem and achieves safe high-speed racing by introducing a progress variable and maximizing track progress. Experiments in simulation and real-world environments validate the proposed method, showing that it can significantly improve the safety of the racing car while maintaining high racing performance.
    \item For vision-based racing control on previously unseen tracks, this thesis proposes a two-stage reinforcement learning-based end-to-end racing method. In the first stage, a teacher policy is trained by reinforcement learning algorithms, which utilizes privileged information in the simulation to achieve high-speed racing control. In the second stage, a student policy that relies solely on visual inputs is trained by knowledge distillation from the teacher policy, enabling the student policy to achieve high-performance racing on unseen tracks. This method improves the training efficiency of vision-based control policies while addressing the challenge of obtaining optimal behaviors in partially observable environments. To enhance the generalization ability of the student policy, domain randomization techniques are also introduced to randomize environment parameters during training. Experimental results show that the proposed method outperforms model-based control methods and other end-to-end methods in racing performance and successfully achieves zero-shot sim-to-real transfer.
    \item For the point-to-point high-speed navigation in unknown complex environments, this thesis proposes an end-to-end high-speed navigation method based on visual predictive models. First, a recurrent state-space model is designed to predict future environmental states. By collecting trajectory data, the dynamic characteristics of the environment are learned from visual inputs, improving the accuracy of state prediction. Based on this visual predictive model, this thesis introduces a model predictive path integral control method that optimizes control strategies by sampling the action space to obtain predicted trajectories. Additionally, for high-speed navigation problems, this thesis designs a reward function to evaluate the quality of sampled trajectories and guide control strategy optimization. This method addresses the sample efficiency problem of reinforcement learning methods and the short-sightedness problem of long-distance navigation in partially observable environments. Finally, comparative and ablation experiments demonstrate the superiority of the proposed method, while real-world experiments further confirm its practical effectiveness.

\end{enumerate}

\noindent \textbf{Keywords}: Autonomous Racing Cars; Model Predictive Control; Control Barrier Function; Reinforcement Learning; End-to-End Navigation